\begin{slikaDesno}{fig/kalibracija_B.pdf}
\PID
Сензор за мерење магнетске инудкције калибрише се помоћу танаког референтног намотаја израђеног од $N = 1000$ 
кружних завојака танке бакарне жице. Полупречник намојатаја је $a = (40\pm 1)\unit{mm}$ а сензор, физички малих димензија, постављен 
је у центру тог прстена. Сензор мери интензитет магнетске инукције $|\bf{B}|$, а његов излаз је напонски, сразмеран том интензитету
према непонзатом коефицијенту $V_{\rm I} = k |\bf{B}|$. Калибрација је обављена мерењем одзива сензора $V_I$ за више различитих 
побудних струја $I$.Апсолутне грешке мерења напона и струје су $\Delta V = 10\unit{\upmu V}$ и $\Delta I = 1\unit{mA}$, а 
резултати мерења су дати у табели \ref{tab:\ID}.
\end{slikaDesno}
\begin{enumerate}[label=(\alph*)]
    \item Одредити непознати калибрациони коефицијент $k$, применом методе најмањих квадрата за линеаризовану теоријску зависност 
    поља од побудне струје.
    \item Проценити апсолутну грешку тог калибрационог коефицијента, $\Delta k$.
\end{enumerate}
\begin{table}[ht!]
    \centering
    \begin{tabular}{c||ccccccccc} 
        \hline
        $I\unit{[mA]}$ & 10 & 20 & 30 & 40 & 50 & 60 & 70 & 80 & 90 \\
        $V_{\rm I}\unit{[mV]}$ & 2,40 & 5,69 & 15,01 & 18,04 & 17,38 & 26,72 & 31,39 & 35,57 & 35,00 \\
        \hline
    \end{tabular}
    \caption{Калибрациона мерења.}
    \label{tab:\ID}
\end{table}

\RESENJE

\begin{table}[b!]
    \centering
    \begin{tabular}{c||ccccccccc} 
        \hline
        $I\unit{[mA]}$ & 10 & 20 & 30 & 40 & 50 & 60 & 70 & 80 & 90 \\
        $V_{\rm I}\unit{[mV]}$ & 2,40 & 5,69 & 15,01 & 18,04 & 17,38 & 26,72 & 31,39 & 35,57 & 35,00 \\ 
        $|{\bf B}|\unit{[\upmu T]}$ & 50,0 & 100 & 150 & 200 & 250 & 300 & 350 & 400 & 450   \\ 
        \hline
    \end{tabular}
    \caption{Допуна табеле.}
    \label{tab:\ID:dopuna}
\end{table}

(а) Магнетска индукција кружне струјне контуре у њеном центру дата је изразом ${\bf B}_0 = \dfrac{\upmu_0 I}{2a} {\bf n}$, где су 
$a$ полупречник контуре, $I$ јачина струје успостављена у контури, ${\bf n}$ вектор нормале везан правилом десне завојнице са 
смером струје, и $\upmu_0 = \dfrac{2\uppi}{5} \unit{\dfrac{\upmu H}{m}}$ је егзактна вредност магнетске пермеабилности вакуума.
Пошто је намотај израђен од $N$ завојака жице, индукција појединачних завојака се сабира (суперпозиција), па је 
${\bf B} = N {\bf B_0} = \dfrac{\upmu_0 N I}{2a} {\bf n}$. Интензитет индукције се онда може одредити на основу 
стује која побудуђује намотај, 
$|{\bf B}| = \dfrac{\upmu_0 N I}{2a}$, и може се израчунати за свако појединачно мерење, па се табела 
\eqref{tab:\ID} може допунити редовима за магнетску индукцију чиме се добија Табела \ref{tab:\ID:dopuna}.

Пошто по претпоставци задатка постоји линеарна веза између напона и поља, онда се вредност коефицијента пропорционалности $k$ 
може потражити помоћу методе најмањих квадрата. Према резултату задатка \ref{z:mnk_1}, односно према изразу 
\eqref{eq:mnk_1:rezultat}, се има
\begin{equation}
    k = \dfrac{\displaystyle
        \sum_{i = 1}^n |{\bf B}|_i V_i
     }{\displaystyle
        \sum_{i = 1}^n V_i^2
     }
     = 11,78 \unit{\dfrac{\upmu T}{
        mV}}
\end{equation}


%Грешка овако одређеног поља онда се 
%може одредити, на основу резултата задатка \ref{z:proizvod} као $\updelta B = \updelta I + \updelta a$